\section{Deterministic single-tape Turing machines}
\label{sec:turing-machines-01}

\begin{definition}[Turing machine]
A \emph{deterministic single-tape Turing machine} is a tuple
\[
  \mathcal T=(Q,\mathcal A_{\mathrm{in}},\mathcal A_{\mathrm{tape}},
  \delta,q_0,q_{\mathrm{acc}},q_{\mathrm{rej}})
\]
such that:
\begin{enumerate}
  \item \(Q\) is a finite set of states;
  \item \(\mathcal A_{\mathrm{in}}\) is a finite input alphabet not containing
        the blank symbol \(\sqcup\);
  \item \(\mathcal A_{\mathrm{tape}}\) is a finite tape alphabet with
        \(\mathcal A_{\mathrm{in}}\cup\{\sqcup\}
        \subseteq\mathcal A_{\mathrm{tape}}\);
  \item \(q_0,q_{\mathrm{acc}},q_{\mathrm{rej}}\in Q\), with
        \(q_{\mathrm{acc}}\neq q_{\mathrm{rej}}\); and
  \item the transition function is
  \[
    \delta:(Q\setminus\{q_{\mathrm{acc}},q_{\mathrm{rej}}\})
    \times\mathcal A_{\mathrm{tape}}
    \longrightarrow Q\times\mathcal A_{\mathrm{tape}}\times\{L,R\}.
  \]
\end{enumerate}
The two \(\mathcal A\)-symbols name components of this machine tuple only;
they are not formula contexts or arithmetical-hierarchy class names.
The tape is infinite in both directions, has only finitely many nonblank cells
at every finite time, and the head scans one cell.
\end{definition}

\begin{definition}[Configuration and computation]
A \emph{configuration} of \(\mathcal T\) is a tuple
\[
  C=(q,T,h),
\]
where \(q\in Q\), \(h\in\Z\), and
\(T:\Z\to\mathcal A_{\mathrm{tape}}\) has finite nonblank support:
\[
  \set{i\in\Z:T(i)\ne\sqcup}
  \quad\text{is finite}.
\]
Suppose \(q\notin\set{q_{\mathrm{acc}},q_{\mathrm{rej}}}\) and
\[
  \delta(q,T(h))=(q',a,d).
\]
Let \(T':\Z\to\mathcal A_{\mathrm{tape}}\) be given by
\[
  T'(i)=
  \begin{cases}
    a,&i=h,\\
    T(i),&i\ne h,
  \end{cases}
\]
and let \(h'=h-1\) when \(d=L\), while \(h'=h+1\) when \(d=R\).  Then, and
only then,
\[
  (q,T,h)\longrightarrow_{\mathcal T}(q',T',h').
\]
A configuration in either halting state has no successor.  Write
\(\longrightarrow_{\mathcal T}^{*}\) for the reflexive transitive closure of
this one-step relation.

For an input word \(w\in\mathcal A_{\mathrm{in}}^*\) of length \(n\), define
\[
  T_w(i)=
  \begin{cases}
    w(i),&0\leq i<n,\\
    \sqcup,&\text{otherwise}.
  \end{cases}
\]
The initial configuration is \((q_0,T_w,0)\).  Thus on empty input the head is
at cell \(0\), which contains a blank.  A computation on \(w\) is the unique
maximal finite or infinite sequence beginning with this initial configuration
and following \(\longrightarrow_{\mathcal T}\) at each successor step.  It
\emph{halts} iff the sequence reaches either halting state, and it
\emph{accepts} or \emph{rejects} according to which halting state it reaches.
\end{definition}

Fix the canonical binary word \(\operatorname{bin}(n)\) for each
\(n\in\N\): \(\operatorname{bin}(0)\) is the one-symbol word \(0\), and for
\(n>0\) it is the ordinary base-two expansion without leading zeros, with the
most significant bit first.  Use the numerical pairing map fixed in the
finite-string chapter:
\[
  \operatorname{pair}_{\N}(a,b)
  =\frac{(a+b)(a+b+1)}2+b.
\]
Its arguments and value are natural numbers; it is distinct in notation and
type from the set-theoretic ordered pair \(\tuple{a,b}\).  For each fixed
arity \(k\), define an injective tuple code by
\[
 \operatorname{code}_0()\defeq0,
 \qquad
 \operatorname{code}_1(n)\defeq n,
\]
and, for \(k\geq2\),
\[
 \operatorname{code}_k(n_0,\ldots,n_{k-1})
 \defeq
 \operatorname{pair}_{\N}\bigl(n_0,
   \operatorname{pair}_{\N}(n_1,\ldots,
     \operatorname{pair}_{\N}(n_{k-2},n_{k-1})\ldots)\bigr).
\]
The arity is fixed by the type \(\N^k\), so it need not be stored inside the
code.  A numerical input for a \(k\)-ary function is the word
\(\operatorname{bin}(\operatorname{code}_k(\bar n))\), placed with its most
significant bit in cell \(0\).

For machines used to compute numerical functions, the input alphabet contains
the symbols \(0,1\).  If a halting configuration has tape function \(T\), let
\[
  m\defeq\min\set{i\in\N:T(i)=\sqcup}.
\]
Finite nonblank support makes \(m\) defined.  The tape prefix of length \(m\)
is a \emph{valid numerical output} iff it is exactly
\(\operatorname{bin}(y)\) for a unique \(y\in\N\).  Thus the empty word, a
word containing a symbol other than \(0,1\), and a nonzero binary word with a
leading zero are invalid.  Cells outside this prefix and the final head
position do not affect the output.  A machine used to compute a numerical
function is required to halt in \(q_{\mathrm{acc}}\) with a valid numerical
output.

\begin{definition}[Partial computable function]
\label{def:partial-computable}
A partial function \(f:\N^k\rightharpoonup\N\) is \emph{partial computable}
iff some Turing machine satisfies both of the following conditions for every
\(\bar n\in\N^k\): when \(f(\bar n)\) is defined, its computation on
\(\operatorname{bin}(\operatorname{code}_k(\bar n))\) reaches
\(q_{\mathrm{acc}}\) with output word \(\operatorname{bin}(f(\bar n))\); and
when \(f(\bar n)\) is undefined, that computation is infinite.  A total
partial computable function is called \emph{computable}.
\end{definition}

The arrow \(\rightharpoonup\) permits undefined values.  Divergence is not a
third numerical output.

The fixed binary and tuple encodings have computable encoding and decoding
maps.  For pairing, given \(z\), boundedly search for the greatest \(s\leq z\)
with \(s(s+1)/2\leq z\), put
\(b=z-s(s+1)/2\), and put \(a=s-b\).  When the arity \(k\) is supplied,
iterating these projections decodes the \(k\) tuple entries; the code alone
does not encode \(k\).  Repeated division by two produces the binary
word, and the left-to-right recurrence \(r\mapsto2r+b\) decodes it.  Each
description is a terminating Turing-machine procedure, so the uses of
``code'' in the definition above do not add an unexamined computation model.

\begin{example}[Unary successor]
A machine computing \(S(n)=n+1\) reads the arity-one input word
\(\operatorname{bin}(n)\), since \(\operatorname{code}_1(n)=n\).  On that block it
scans right to the end, propagates a carry left through consecutive \(1\)s,
 and writes \(1\) at the first \(0\) encountered.  If the carry passes the most
 significant bit of an \(m\)-bit input, all old bits are now \(0\).  The
 machine shifts that \(m\)-cell zero block one cell to the right, writes \(1\)
 in cell \(0\), and leaves a blank in cell \(m+1\).  Thus the output prefix in
 cells \(0,\ldots,m\) changes \(111\) to \(1000\).  It then treats that prefix as the output
word, not as another tuple code, because the output convention requires the
binary word for the numerical value.  The loop invariant is that every
processed suffix is the correct low-order suffix of \(n+1\), while the state
records whether a carry remains.  Decoding and block relocation are justified
by the computable coding operations established above.
\end{example}

\begin{warning}[Machine formalism versus complexity measure]
Changing from one tape to several tapes does not change which partial functions
are computable, but it can change running time by a polynomial factor.  A
computability claim and a complexity claim therefore require different
invariance arguments; see \cite{BoolosBurgessJeffrey2007} for the
multi-tape invariance argument.
\end{warning}

